%%%%%%%%%%%%%%%%%%%%%%%%%%%%%%%%%%%%%%%%%
% Short Sectioned Assignment LaTeX Template Version 1.0 (5/5/12)
% This template has been downloaded from: http://www.LaTeXTemplates.com
% Original author:  Frits Wenneker (http://www.howtotex.com)
% License: CC BY-NC-SA 3.0 (http://creativecommons.org/licenses/by-nc-sa/3.0/)
%%%%%%%%%%%%%%%%%%%%%%%%%%%%%%%%%%%%%%%%%

% \documentclass[paper=a4, fontsize=11pt]{scrartcl} % A4 paper and 11pt font size
\documentclass[11pt, a4paper]{book}
\usepackage[T1]{fontenc} % Use 8-bit encoding that has 256 glyphs
\usepackage[utf8]{inputenc}
\usepackage{fourier} % Use the Adobe Utopia font for the document - comment this line to return to the LaTeX default
\usepackage{listings} % para insertar código con formato similar al editor
\usepackage[spanish, es-tabla]{babel} % Selecciona el español para palabras introducidas automáticamente, p.ej. "septiembre" en la fecha y especifica que se use la palabra Tabla en vez de Cuadro
\usepackage{url} % ,href} %para incluir URLs e hipervínculos dentro del texto (aunque hay que instalar href)
\usepackage{graphics,graphicx, float} %para incluir imágenes y colocarlas
\usepackage[gen]{eurosym} %para incluir el símbolo del euro
\usepackage{cite} %para incluir citas del archivo <nombre>.bib
\usepackage{enumerate}
\usepackage{hyperref}
\usepackage{graphicx}
\usepackage{tabularx}
\usepackage{booktabs}

\usepackage[table,xcdraw]{xcolor}
\hypersetup{
	colorlinks=true,	% false: boxed links; true: colored links
	linkcolor=black,	% color of internal links
	urlcolor=cyan		% color of external links
}
\renewcommand{\familydefault}{\sfdefault}
\usepackage{fancyhdr} % Custom headers and footers
\pagestyle{fancyplain} % Makes all pages in the document conform to the custom headers and footers
\fancyhead[L]{} % Empty left header
\fancyhead[C]{} % Empty center header
% \fancyhead[R]{Me Llamo Así} % My name
\fancyfoot[L]{} % Empty left footer
\fancyfoot[C]{} % Empty center footer
\fancyfoot[R]{\thepage} % Page numbering for right footer
%\renewcommand{\headrulewidth}{0pt} % Remove header underlines
\renewcommand{\footrulewidth}{0pt} % Remove footer underlines
\setlength{\headheight}{13.6pt} % Customize the height of the header

\usepackage{titlesec, blindtext, color}
\definecolor{gray75}{gray}{0.75}
\newcommand{\hsp}{\hspace{20pt}}
\titleformat{\chapter}[hang]{\Huge\bfseries}{\thechapter\hsp\textcolor{gray75}{|}\hsp}{0pt}{\Huge\bfseries}
\setcounter{secnumdepth}{4}
\usepackage{sectsty}
\usepackage{titlesec}

\renewcommand{\chaptermark}[1]{%
  \markboth{Capítulo \thechapter.\ #1}{}}

\renewcommand{\sectionmark}[1]{%
  \markright{\thesection.\ #1}}


\begin{document}

	% Plantilla portada UGR
	\begin{titlepage}
 
 
\newlength{\centeroffset}
\setlength{\centeroffset}{-0.5\oddsidemargin}
\addtolength{\centeroffset}{0.5\evensidemargin}
\thispagestyle{empty}

\noindent\hspace*{\centeroffset}\begin{minipage}{\textwidth}

\centering
\includegraphics[width=0.9\textwidth]{imagenes/logo.png}\\[1.4cm]

\textsc{ \Large TRABAJO FIN DE GRADO\\[0.2cm]}
\textsc{ GRADO EN INGENIERÍA INFORMÁTICA}\\[1cm]
% Upper part of the page
% 
% Title
{\Huge\bfseries \\ Desarrollo de una aplicaci´on
de modelado 3D}
\noindent\rule[-1ex]{\textwidth}{3pt}\\[3.5ex]
{\large\bfseries Motor de renderizado en tiempo real basada en SDFs}
\end{minipage}

\vspace{2.5cm}
\noindent\hspace*{\centeroffset}\begin{minipage}{\textwidth}
\centering

\textbf{Autor}\\ {Pablo Cantudo Gómez}\\[2.5ex]
\textbf{Directores}\\
{Juan Carlos Torres Cantero y Luis López Escudero\\}\\[2cm]
\includegraphics[width=0.3\textwidth]{imagenes/etsiit_logo.png}\\[0.1cm]
\textsc{Escuela Técnica Superior de Ingenierías Informática y de Telecomunicación}\\
\textsc{---}\\
Granada, Septiembre de 2025
\end{minipage}
%\addtolength{\textwidth}{\centeroffset}
%\vspace{\stretch{2}}
\end{titlepage}




	% Plantilla prefacio UGR
	\thispagestyle{empty}

\begin{center}
{\large\bfseries Desarrollo software de una aplicación de modelado 3d basada en Signed Distance Funtion (sdf) \\ Motor de renderizado en tiempo real usando Dawn, C++ y WebGPU }\\
\end{center}
\begin{center}
Pablo Cantudo Gómez\\
\end{center}

\vspace{0.7cm}

\vspace{0.5cm}
\noindent\textbf{Palabras clave}: \textit{SDF}, \textit{WebGPU}, \textit{C++}, \textit{ray marching}, \textit{motor de renderizado}
\vspace{0.7cm}

\noindent\textbf{Resumen}\\
\break
Este trabajo presenta el desarrollo de \textit{Copper}, un motor de renderizado 3D en tiempo real implementado en C++ utilizando la API WebGPU a trav\'es de Dawn.
El motor emplea t\'ecnicas de \textit{ray marching} sobre funciones de distancia con signo (SDF), lo que permite una representaci\'on eficiente y flexible de geometr\'ia 
 tridimensional. Una de las principales ventajas de utilizar SDF es la capacidad de expresar de forma directa y compacta operaciones booleanas entre primitivas geom\'etricas
  ---como uniones, intersecciones o diferencias--- mediante simples expresiones algebraicas. Adem\'as, a diferencia de los m\'etodos tradicionales basados en mallas o 
  pol\'igonos, las SDF permiten aplicar operaciones booleanas suaves, como la uni\'on con suavizado (\textit{smooth union}), 
  de forma pr\'acticamente trivial desde el punto de vista computacional.
\bigbreak
Este enfoque simplifica notablemente la construcci\'on de formas complejas, evitando problemas t\'ipicos de la computaci\'on geom\'etrica como el manejo 
de v\'ertices, normales o topolog\'ias complejas. Tambi\'en se facilita la animaci\'on y transformaci\'on de objetos mediante funciones continuas. El motor 
incluye un sistema de sombreado en WGSL, con soporte para sombras suaves, operaciones modulares sobre la escena, seleccion de objetos, carga y guardado de modelos. Los resultados demuestran que
el uso de SDF no solo ofrece un modelo m\'as elegante para definir geometr\'ia, sino que permite construir escenas visualmente complejas con menos c\'odigo y una mayor expresividad gr\'afica.
En conjunto, el sistema demuestra la viabilidad t\'ecnica y creativa del uso de SDF en entornos modernos.
\cleardoublepage

\begin{center}
	{\large\bfseries Development of a 3D modeling application based on Signed Distance Functions (SDF) \\ Real-time rendering engine using Dawn, C++ and WebGPU}\\
\end{center}
\begin{center}
	Pablo Cantudo Gómez\\
\end{center}
\vspace{0.5cm}
\noindent\textbf{Keywords}: \textit{SDF}, \textit{WebGPU}, \textit{C++}, \textit{ray marching}, \textit{rendering engine}
\vspace{0.7cm}

\noindent\textbf{Abstract}\\
\break
This work presents the development of Copper, a real-time 3D rendering engine implemented in C++ using the WebGPU API through Dawn.
The engine employs ray marching techniques over Signed Distance Functions/Fields (SDF), enabling an efficient and flexible representation of three-dimensional geometry.
One of the main advantages of using SDF is the ability to directly and compactly express Boolean operations between geometric primitives ---such as unions, intersections, or differences--- through simple algebraic expressions.
Furthermore, unlike traditional mesh- or polygon-based methods, SDF allows for smooth Boolean operations, such as smooth union, in a virtually trivial manner from a computational standpoint.
\bigbreak
This approach greatly simplifies the construction of complex shapes, avoiding typical problems in geometric computing such as handling vertices, normals, or complex topologies.
It also facilitates the animation and transformation of objects through continuous functions.
The engine includes a shading system in WGSL, with support for soft shadows, modular scene operations, object selection, and model loading and saving.
The results show that the use of SDF not only offers a more elegant model for defining geometry, but also enables the creation of visually complex scenes with less code and greater graphical expressiveness.
Overall, the system demonstrates the technical and creative feasibility of using SDF in modern environments.

\cleardoublepage

\thispagestyle{empty}

\noindent\rule[-1ex]{\textwidth}{2pt}\\[4.5ex]

D. \textbf{Tutora/e(s)}, Profesor(a) del ...

\vspace{0.5cm}

\textbf{Informo:}

\vspace{0.5cm}

Que el presente trabajo, titulado \textit{\textbf{Chief}},
ha sido realizado bajo mi supervisión por \textbf{Estudiante}, y autorizo la defensa de dicho trabajo ante el tribunal
que corresponda.

\vspace{0.5cm}

Y para que conste, expiden y firman el presente informe en Granada a Junio de 2018.

\vspace{1cm}

\textbf{El/la director(a)/es: }

\vspace{5cm}

\noindent \textbf{(nombre completo tutor/a/es)}

\chapter*{Agradecimientos}
Quiero expresar mi agradecimiento a todas las personas que han contribuido directa o indirectamente a la realización de este trabajo.\break

En primer lugar, a mi tutor Juan Carlos Torres Cantero y cotutor Luis López Escudero por la ayuda durante el desarrollo y la correci\'on del proyecto,
a mis amigos por las discusiones y las ideas, y a mi familia por su apoyo incondicional.\break

 Finalmente, me gustaría mencionar a la comunidad de desarrolladores y recursos abiertos, especialmente los foros, artículos y proyectos relacionados con WebGPU, SDF y ray marching,
 cuya documentación y ejemplos han sido una fuente de aprendizaje invaluable.





	% Índice de contenidos
	\newpage
	\tableofcontents

	% Índice de imágenes y tablas
	\newpage
	\listoffigures

	% Si hay suficientes se incluirá dicho índice
	\listoftables 
	\newpage

	% Introducción 
	\chapter{Introducción}


	% Descripción del problema y hasta donde se llega
	\chapter{Descripción del problema}



	% Estado del arte
	% 	1. Crítica al estado del arte
	% 	2. Propuesta
	\chapter{Estado del arte}

	
	\chapter{Planificación}

\section{Metodología utilizada}


\section{Temporización}

\section{Seguimiento del desarrollo}


	% Análisis del problema
	% 1. Análisis de requisitos
	% 2. Análisis de las soluciones
	% 3. Solucion propuesta
	% 4. Análisis de seguridad
	\chapter{Análisis del problema}
 


	% Desarrollo bajo sprints: 
	% 	1. Permitir registros y login de usuarios
	% 	2. Desarrollo del sistema de incidencias
	% 	3. Desarrollo del sistema de denuncias administrativas y accidentes
	% 	4. Desarrollo del sistema de croquis
	%   5. Instalación de la aplicación de manera automática
	\chapter{Implementación}

La implementación del software se ha dividido en hitos. Estos, han sido definidos en Github
y cada uno de ellos contiene un grupo de \textit{issues} que se corresponden con las distintas
mejoras que se han ido incorporando al software a lo largo de su desarrollo.\\



	% Presupuesto

	% Conclusiones
	\chapter{Conclusiones y trabajos futuros}



	% Trabajos futuros


	
	\newpage
	\bibliography{bibliografia}
	\bibliographystyle{plain}
	
\end{document}

